% Generated draft: requires mathematical and source review.
\documentclass[11pt]{article}
\usepackage[margin=25mm]{geometry}
\usepackage{fontspec}
\setmainfont{Latin Modern Roman}
\usepackage{amsmath,amssymb,mathtools,unicode-math}
\setmathfont{Latin Modern Math}
\usepackage{enumitem,needspace,xurl,hyperref}
\setlength{\parindent}{0pt}
\setlength{\parskip}{5pt}
\setlength{\emergencystretch}{3em}
\newcommand{\N}{\mathbb{N}}
\newcommand{\Z}{\mathbb{Z}}
\newcommand{\Q}{\mathbb{Q}}
\newcommand{\R}{\mathbb{R}}
\newcommand{\C}{\mathbb{C}}
\newcommand{\F}{\mathbb{F}}
\newcommand{\A}{\mathbb{A}}
\newcommand{\PP}{\mathbb{P}}
\newcommand{\E}{\mathbb{E}}
\newcommand{\eps}{\varepsilon}
\newcommand{\dd}{\,\mathrm{d}}
\begin{document}
\section*{UnsolvedMath ID: 800211}
\textbf{Source problem code:} catalogue entry
\subsection{Definitions and mathematical statement}
Let \(D\subset\C^2\) be a strictly pseudoconvex domain diffeomorphic, as a real smooth manifold, to \(\R^4\). Strict pseudoconvexity is understood in the usual boundary sense: locally a defining function \(\rho\) has positive Levi form on nonzero complex tangent vectors to \(\partial D\). The source does not specify boundary regularity; this definition requires at least a \(C^2\) defining function. No boundedness assumption is supplied.

Write \(\overline\Delta=\{\zeta\in\C:|\zeta|\leq1\}\). An embedded analytic disc \(A\) is the image of an embedding of \(\overline\Delta\) whose restriction to \(\Delta\) is holomorphic; the required boundary regularity of this embedding is not specified in the source. Write \(A^\circ\) for its interior and \(bA\) for its boundary. Assume that \(A\) lies ``outside \(D\)'' and that
\[
 bA\subset\partial D.
\]
The exterior condition entails \(A^\circ\cap D=\varnothing\), but the source does not say whether it additionally requires \(A^\circ\cap\overline D=\varnothing\).

The source first asserts the existence of a smoothly embedded closed real disc \(M\) with
\[
 M^\circ\subset D,\qquad bM=bA,
\]
which is totally real except at one elliptic complex tangent point \(p\in M\). Here total reality at \(q\) means
\[
 T_qM\cap iT_qM=\{0\}.
\]
An elliptic point is understood in the usual Bishop sense: in suitable local holomorphic coordinates centered at \(p\), the surface has the form
\[
 w=z\overline z+\lambda(z^2+\overline z^{\,2})+o(|z|^2),
 \qquad 0\leq\lambda<\tfrac12.
\]
The interior containment above interprets the source's phrase ``\(M\) in \(D\)'' consistently with its boundary condition; it does not put \(bM\) inside the open domain \(D\).

\textbf{Question.} For such a configuration \((M,A)\subset\C^2\), must the envelope of holomorphy of \(M\) contain \(A\)? Symbolically, with the unresolved meaning of \(\mathcal E(M)\) retained,
\[
 A\subset\mathcal E(M)\ ?
\]
An envelope of holomorphy is ordinarily defined for an open domain or a Riemann domain, whereas \(M\) here is a real disc. The source supplies no definition using neighborhoods or holomorphic germs along \(M\), and does not specify how containment of \(A\) is to be interpreted if the envelope has multiple sheets. Thus the displayed question retains an unresolved construction, rather than defining a precise envelope. The wording also does not explicitly distinguish a conclusion for every eligible choice of \(M\) from the existence of a suitable choice of \(M\). The preceding existence assertion is recorded as part of the supplied formulation, without an independent proof or verification.
\subsection{Short English statement}
The source asks whether an exterior embedded analytic disc attached to the boundary of a strictly pseudoconvex domain must lie in the envelope of holomorphy of an interior smooth spanning disc having exactly one elliptic complex tangent. The meaning of the envelope for this real disc, the exterior condition, and the quantifier over spanning discs require clarification. The supplied registry labels the problem open; that status is not independently verified.
\subsection{Sources}
\begin{itemize}
\item Supplied local UnsolvedMath registry record: permanent ID \(800211\), exact source problem code \texttt{catalogue entry}, title \emph{Several Complex Variables Problem List --- Problem 7}. The supplied status is \texttt{open}.
\item Buzzard, Fornaess, Krantz, \emph{The Several Complex Variables Problem List} (1995), Problem 7. Supplied source URL: \url{https://arxiv.org/abs/math/9509202}. The registry's proposer field names Buzzard, Fornaess, Krantz; the quoted problem itself is attributed to F. Forstneric, with the supplied contact \texttt{forstner@math.wisc.edu}.
\item Supplied provenance: accessed 2026-10-16; extraction method \texttt{arxiv\_latex}; upstream revision \texttt{d371ee917ac9647fbb06cd6fc529deec6f3dda82}. This access date and provenance are supplied curation information, not evidence of independent source consultation or verification in this draft.
\end{itemize}
\end{document}
